\documentclass[11pt]{report}

% --- Basics ---
\usepackage[a4paper,margin=1in]{geometry}
\usepackage[hidelinks]{hyperref}
\usepackage{graphicx}
\usepackage{float}      % for [H] figure placement when needed
\usepackage{caption}
\usepackage{booktabs}   % nice tables (optional)

% Optional: if you want better crossrefs like "Figure 3.2"
% \usepackage[nameinlink,noabbrev]{cleveref}

% Optional: captions for many screenshots can get repetitive
% \captionsetup[figure]{font=small}

% --- Screenshot helpers (put in the preamble) ---
\usepackage{graphicx}
\usepackage{float}      % enables [H]
\usepackage{caption}

% Optional: common folder for images
\newcommand{\imgdir}{figs/}

% Main screenshot macro:
%  - #1: filename (relative to \imgdir unless you include a path)
%  - #2: caption text
%  - #3: label (without "fig:"; the macro adds it)
%  - #4: width (optional; default 0.98\linewidth)
\newcommand{\screenshot}[4][0.98\linewidth]{%
    \begin{figure}[H]
        \centering
        \includegraphics[width=#1]{\imgdir#2}
        \caption{#3}
        \label{fig:#4}
    \end{figure}
}

% Variant: screenshot with a short caption for LoF + long caption for the page
%  - #1 width (optional; default 0.98\linewidth)
%  - #2 filename
%  - #3 short caption (List of Figures)
%  - #4 long caption (printed under figure)
%  - #5 label key
\newcommand{\screenshotLF}[5][0.98\linewidth]{%
    \begin{figure}[H]
        \centering
        \includegraphics[width=#1]{\imgdir#2}
        \caption[#3]{#4}
        \label{fig:#5}
    \end{figure}
}

% Quick reference helper for figures
\newcommand{\figref}[1]{Figure~\ref{fig:#1}}


\title{Transgraphier\\{\large User Manual}}
\author{Fernando Cacciola}
\date{Version 1.0 --- \today}

\begin{document}
    \maketitle
    \tableofcontents
    
    % -------------------------
    \chapter*{About This Manual}
    \addcontentsline{toc}{chapter}{About This Manual}
    \section*{Version, Audience, and Conventions}
    \addcontentsline{toc}{section}{Version, Audience, and Conventions}
    % - Which app version this matches
    % - Intended audience (end-users / analysts)
    % - Notation, terminology, UI conventions (menu names, buttons, etc.)
    
    \section*{Related Documents}
%    \addcontentsline{toc}{section}{Related Documents}
%    You can read the following document in order to understand what
%    is it that this Application does at a technical level:
%    \href{run:TechnicalPaper.pdf}{Technical Paper (PDF)}.
    
    % -------------------------
    \chapter{Introduction}
    \section{What the App Does}
       
       Given an \textbf{EVP} audio file that is \textit{specifically} is expected
       to contain a text message \textit{encoded in binary} (like computers do),
       but using \textbf{Tap Code} as encoding system, applies a sequence of audio processing steps in order to attempt to extract a Text Message from it.
       
    \section{Typical Workflow at a Glance}
    
\begin{itemize}
    \item Open an audio file containing an \textbf{EVP}
         
    \item Click on Process (or Re-Process if this one file has already been processed)
         
    \item Look at the \textbf{TOP-RIGHT} panel.\\
          If will either say that ``\textit{no message was decoded}'' OR\\
          It will contain the text message that was decoded from the \textbf{EVP} audio.
         
    \item The TOP-LEFT panel will show a general log detailing what\\
          the processing was doing.
         
    \item You can also look at the sequence of ``Filter Outputs'', right below\\
           the Input Wave Form.\\
           Details about the result of each processing step are shown there.
\end{itemize}
    
    
    % Open WAV -> Process -> Review slots + waveforms -> Export/Save (if relevant)
    \section{Key Concepts and Terminology}
    % "slot", "segment", "window", "analysis result", etc.
    
    % -------------------------
    \chapter{Installation and Setup}
    \section{System Requirements}
    \section{Installation}
    \section{Updating / Uninstalling}
    \section{First Launch Checklist}
    % where files are stored, default folders, permissions, audio formats supported
    
    % -------------------------
    \chapter{Quick Start (5 Minutes)}
    \section{Open a WAV File}
    \section{Run \texttt{Process}}
    \section{Read the Results View}
    % interleaved text slots + waveforms explanation at a very high level
    \section{Common First Issues}
    % "no output", "file not supported", "long processing time", etc.
    
    % -------------------------
    \chapter{User Interface Tour}
    \section{Main Window Layout}
    
    % screenshot: full window labeled
    \section{Menus and Toolbar}
    \subsection{File Menu}
    \subsection{Process Menu}
    \subsection{View Menu}
    \subsection{Help Menu}
    \section{Panels / Regions}
    \subsection{Input Waveform View}
    \subsection{Results Timeline / Slots View}
    \subsection{Details / Inspector (if any)}
    \subsection{Status / Log Area (if any)}
    
    % -------------------------
    \chapter{Primary Workflow: From WAV to Results}
    \section{Selecting and Loading Audio}
    \subsection{Supported Formats and Constraints}
    \subsection{Large Files and Performance Notes}
    \section{Running Processing}
    \subsection{Process Options (High-Level)}
    \subsection{Progress, Cancellation, and Errors}
    \section{Understanding the Output}
    \subsection{What a Text Slot Represents}
    \subsection{What Each Waveform Represents}
    \subsection{Navigation (zoom, pan, selection)}
    \subsection{Interpreting Common Patterns}
    % include representative screenshots here
    \section{Saving / Exporting Results (if applicable)}
    \subsection{Export Formats}
    \subsection{Reproducibility (saving presets, parameters)}
    
    % -------------------------
    \chapter{Working with Results}
    \section{Browsing and Filtering Slots}
    \section{Comparing Regions / Slots}
    \section{Measurements and Annotations (if applicable)}
    \section{Copying / Sharing Outputs}
    % images, CSV, reports, etc.
    
    % -------------------------
    \chapter{Processing Options and Presets}
    \section{Parameters Overview}
    % keep user-facing; link deep math to technical paper
    \section{Presets}
    \subsection{Choosing a Preset}
    \subsection{Creating / Editing Presets}
    \subsection{Reset to Defaults}
    \section{Quality vs Speed Tradeoffs}
    
    % -------------------------
    \chapter{Feature Reference}
    % Use one section per major feature area. This is the "reference" part.
    \section{Audio Input Handling}
    \section{Segmentation / Windowing Controls (if applicable)}
    \section{Detection / Classification Outputs (if applicable)}
    \section{Visualization Controls}
    \subsection{Waveform Display Options}
    \subsection{Markers, Overlays, and Colors (if applicable)}
    \section{Batch Processing (if applicable)}
    \section{Command-line / Automation (if applicable)}
    
    % -------------------------
    \chapter{Troubleshooting}
    \section{Common Symptoms and Fixes}
    \subsection{The app won’t open / crashes}
    \subsection{WAV won’t load / silent audio}
    \subsection{Process produces empty output}
    \subsection{Processing is too slow}
    \subsection{Waveforms look incorrect}
    \section{Error Messages}
    % copy/paste friendly list; map each to resolution steps
    \section{How to Collect Logs for Support}
    
    % -------------------------
    \chapter{FAQ}
    \section{How do I interpret [common output]?}
    \section{How do I choose parameters?}
    \section{Can I reproduce results exactly?}
    \section{What does the algorithm do? (link)}
    See:
    \href{run:TechnicalPaper.pdf}{Technical Paper (PDF)}.
    
    % -------------------------
    \chapter{Keyboard Shortcuts}
    % table
    
    \chapter{Glossary}
    % definitions: slot, segment, frame, envelope, etc.
    
    % -------------------------
    \appendix
    \chapter{File Formats}
    \section{Input Audio Requirements}
    \section{Exported Data Formats}
    
    \chapter{Advanced Notes}
    \section{Algorithm Overview (Non-mathematical)}
    % Keep this gentle; point to paper for proofs/details
    For full details, see:
    \href{run:TechnicalPaper.pdf}{Technical Paper (PDF)}.
    
    \chapter{Change Log}
    % optional
    
\end{document}
